%% data.tex — Section 3: Data and Sample Characteristics

\subsection{State Monitoring of Quality of Primary Education}
\label{subsec:data_monitoring}

I use two waves of Ukraine's State Monitoring of Quality of Primary Education
(hereafter, the Monitoring), a nationally representative assessment programme
that tracks achievement trends, evaluates alignment with national learning
standards, and examines factors associated with student outcomes. The programme
assesses fourth-grade students in mathematics and reading, and collects parallel
questionnaires from students and teachers.

The 2018 baseline sample was drawn from all Ukrainian-language general secondary
schools with fourth-grade enrolment. Schools were stratified by size and selected
via prob\-ability-pro\-por\-tional-to-size (PPS) sampling, yielding 366 schools and
486 classes. Each class was randomly assigned to complete either the mathematics
or the reading assessment, not both. Fieldwork ran from 17~April to 18~May~2018
and produced data on 9,077 tested pupils (4,535 in reading; 4,542 in mathematics);
after standard exclusions for missing data, the analytic sample comprises 9,007
valid records. All 486 sampled teachers completed questionnaires.

The 2021 wave followed identical sampling and administration procedures. The
target sample included 365 schools and 485 classes, but pandemic-related logistical
disruptions reduced the achieved sample to 355 schools and 475 classes. Fieldwork
ran from 16~April to 17~May~2021, yielding 7,991 tested pupils (4,192 in reading;
3,799 in mathematics), of which 7,924 records are retained after exclusions.
A total of 473 teachers completed questionnaires.

In both waves, external proctors administered the assessments under standardised
protocols with secure handling of materials. Student achievement was scaled using
item response theory (IRT): dichotomous items followed a two-parameter logistic
(2PL) model and polytomous items a graded-response model. All responses were placed
on a common logit scale that ensures comparability across rotated test forms and
between waves. Logit scores were linearly transformed to a 100--300-point reporting
scale, with cut-scores at 170, 200, and 230 defining four performance levels
(below-basic, basic, middle, and high). Because both waves use the same IRT
parameters and reporting transformation, scores are directly comparable at the
population level \citep{uceqa2025}. All analyses use survey sampling weights
and cluster standard errors at the school level (primary sampling unit).

\subsection{Sample Comparability and Attrition}
\label{subsec:attrition}

Because the design compares repeated cross-sections rather than the same students
or schools across time, sample comparability across waves is essential.
Table~\ref{tab:balance} presents balance statistics for the mathematics sample.

First, enrolment fell by roughly 550 students (11\%), and the share of enrolled
students who sat the test dropped from 90.8 to 85.8~percent, representing
approximately 720 fewer tested pupils in mathematics. If lower-performing or
more disadvantaged students disproportionately missed the 2021 assessment, the
observed cross-cohort gap understates the true learning loss; if pandemic-disrupted
families were the main absentees, it overstates it. I address this in the
robustness checks.

Second, most demographic variables are well balanced. Gender composition is virtually
identical across waves (SMD~=~0.005), and the rural share increased modestly from
25 to 30~percent (SMD~=~0.099). Age rose by roughly one month (SMD~=~0.161),
consistent with calendar effects or pandemic-related delays in grade progression.

Third, SES composition shifted substantially. The high-SES tertile grew from
21 to 35~percent of tested students, while the low- and middle-SES shares both
declined (combined SMD~=~0.340). By conventional thresholds, an SMD above 0.25
indicates substantial imbalance; at 0.340, the SES shift exceeds this threshold.
Because SES is positively correlated with test scores, a 2021 sample skewed toward
higher-SES students mechanically compresses the cross-cohort gap: the main estimates
are best read as lower bounds on the true learning loss.

Fourth, school-type composition shifted modestly, with gymnasium and lyceum students
rising from 9 to 13~percent of the sample and specialised schools falling from
15 to 12~percent. Books-at-home distributions remained stable (SMD~=~0.145).

The two samples are well matched on demographics but diverge on SES.
The regression analysis controls for all observed differences; the robustness
checks in Section~\ref{sec:robustness} assess sensitivity to the compositional
shift.

\begin{table}[htbp]
\centering
\caption{Sample comparability across waves: Mathematics}
\label{tab:balance}
\begin{tabular}{lccc}
\toprule
\textbf{Variable} & \textbf{2018} & \textbf{2021} & \textbf{SMD} \\
\midrule
\multicolumn{4}{l}{\textit{Sample size}} \\
\quad Enrolled          & 4,956 & 4,409 & \\
\quad Tested            & 4,501 & 3,781 & \\
\quad Proportion tested & 0.908 & 0.858 & \\
\addlinespace
\multicolumn{4}{l}{\textit{Demographics}} \\
\quad Female            & 0.497 & 0.499 & 0.005 \\
\quad Age (months)      & 121.5 & 122.4 & 0.161 \\
\quad Rural             & 0.254 & 0.298 & 0.099 \\
\addlinespace
\multicolumn{4}{l}{\textit{School type (\%)}} \\
\quad Gymnasium/Lyceum  &  8.7 & 12.6 & \\
\quad Regular school    & 76.0 & 75.4 & \\
\quad Specialised       & 15.3 & 12.0 & 0.150 \\
\addlinespace
\multicolumn{4}{l}{\textit{SES tertile (\%)}} \\
\quad Low               & 30.9 & 24.1 & \\
\quad Middle            & 32.2 & 23.4 & \\
\quad High              & 21.2 & 34.5 & \\
\quad Missing           & 15.6 & 18.1 & 0.340 \\
\addlinespace
\multicolumn{4}{l}{\textit{Books at home (\%)}} \\
\quad 0--10             & 12.4 & 17.2 & \\
\quad 11--25            & 35.1 & 35.0 & \\
\quad 26--100           & 31.5 & 29.3 & \\
\quad 101--200          & 11.9 & 10.0 & \\
\quad $>$200            &  7.4 &  7.0 & \\
\quad Missing           &  1.6 &  1.5 & 0.145 \\
\bottomrule
\end{tabular}
\begin{tablenotes}
\small
\item \textit{Notes}: Weighted means or percentages using survey design weights.
SMD~= standardised mean difference. Thresholds: SMD~$>$~0.10 may indicate
imbalance; SMD~$>$~0.25 indicates substantial imbalance.
\end{tablenotes}
\end{table}
